\label{sec:background}

\subsection{Census Bureau Counting Rules and Their History}

The Census Bureau's practice of counting incarcerated people at the prison
location is longstanding, dating to at least the nineteenth century.
The Bureau's position has been that correctional facilities constitute a
form of ``usual residence'' under its criteria, even though incarcerated
people generally lack any choice of location and cannot participate in the
civic life of the prison community. This practice was reaffirmed in the
2010 and 2020 censuses despite growing advocacy for a change in methodology.

In 2020, the Bureau introduced the Detailed Demographic and Housing
Characteristics (DHC-A) file, which provides prison-facility-level data
including total incarcerated population by facility and demographic
breakdown. This file enables states and redistricters to implement
``adjust-at-home'' corrections without relying on Bureau methodology
changes. The National Prisoner Statistics (NPS) program, administered
by the Bureau of Justice Statistics (BJS), provides complementary data
on the home states of federal prisoners.

\subsection{\textit{Evenwel v.\ Abbott} (2016) and Population Base Discretion}

The central Supreme Court precedent governing the population base for
redistricting is \textit{Evenwel v.\ Abbott},\footnote{\textit{Evenwel
v.\ Abbott}, 578 U.S.\ 54 (2016).} decided unanimously in 2016.
\textit{Evenwel} arose from a challenge to Texas Senate districts.
Plaintiffs argued that the Equal Protection Clause required states to
equalize districts based on citizen voting-age population (CVAP) rather
than total population, on the theory that using total population
diluted the votes of citizens in districts with large non-citizen
populations. The Court unanimously rejected this argument.

Writing for the Court, Justice Ginsburg held that states may use total
population as the apportionment base. The majority opinion reviewed the
history of the ``one person, one vote'' doctrine from \textit{Reynolds v.\
Sims}\footnote{\textit{Reynolds v.\ Sims}, 377 U.S.\ 533 (1964).} and
concluded that the Constitution permits---but does not require---states
to equalize total population, voter-eligible population, or any
reasonable alternative. The Court left open whether some alternative
bases might be constitutionally impermissible, a question it explicitly
did not resolve.

Critically, \textit{Evenwel} did not address prison gerrymandering.
The case was about whether states \textit{must} use CVAP rather than total
population; it was not about whether states may reallocate the total
population by counting incarcerated people differently. The \textit{Evenwel}
holding is thus compatible with prison population adjustment: a state that
uses total population but reassigns incarcerated people to their home
counties is still using total population as the base---it is simply
deciding where within total population to count each person.

This distinction is legally important. Prison population adjustment
does not change the apportionment base (total population vs.\ CVAP);
it changes the \textit{location} at which individuals within the total
population base are counted. States retain the discretion to make
this locational assignment under \textit{Evenwel}'s framework,
and at least twelve states have exercised that discretion.

\subsection{Survey of State Prison Population Adjustment Laws}

The following states had enacted prison population adjustment legislation
for redistricting purposes as of 2026. For each state, the key law
and effective date are noted.

\begin{table}[ht]
\centering
\caption{States with Prison Population Adjustment Laws for Redistricting (as of 2026)}
\label{tab:state-laws}
\begin{tabular}{lll}
\toprule
State & Legislation & Year Enacted \\
\midrule
California & Cal.\ Elec.\ Code \S\ 21003 (AB 420) & 2011 \\
Colorado & Colo.\ Rev.\ Stat.\ \S\ 2-1-901 & 2010 \\
Delaware & Del.\ Code Ann.\ tit.\ 29, \S\ 804 & 2010 \\
Illinois & 10 ILCS 5/3A-1 et seq. & 2011 \\
Maryland & Md.\ State Gov't Code \S\ 2-2A-01 & 2010 \\
Michigan & Mich.\ Comp.\ Laws \S\ 168.931 & 2022 \\
Montana & Mont.\ Code Ann.\ \S\ 13-2-109 & 2015 \\
Nevada & Nev.\ Rev.\ Stat.\ \S\ 218C.130 & 2021 \\
New Jersey & N.J. Stat.\ Ann.\ \S\ 52:14-1.4 & 2020 \\
New York & N.Y.\ Legis.\ Law \S\ 83-m & 2010 \\
Virginia & Va.\ Code Ann.\ \S\ 24.2-304.03 & 2020 \\
Washington & Wash.\ Rev.\ Code \S\ 44.05.160 & 2011 \\
\bottomrule
\end{tabular}
\end{table}

The adjustment approach is broadly similar across these states:
incarcerated people are subtracted from the census count of the census
tract in which the prison is located and added to the census count of
their last known home address (usually a county or municipality).
States differ in their data sources for the home-address reallocation,
with some using self-reported data from state corrections departments
and others relying on the Bureau of Justice Statistics NPS data.

\subsection{Racial and Demographic Profile of Incarcerated Populations}

According to the Bureau of Justice Statistics, approximately 58\% of
people held in state and federal correctional facilities are Black or
Hispanic.\footnote{Bureau of Justice Statistics, \textit{Prisoners in
2020--Statistical Tables} (BJS, 2021).} This disproportionality
means that the prison gerrymandering effect falls heavily on communities
of color: their members are counted in rural, predominantly white districts
that contain prisons rather than in the urban and suburban communities
where they lived before incarceration.

This demographic profile gives prison gerrymandering a direct Voting Rights
Act dimension. When Black and Hispanic residents are systematically
undercounted in urban minority-opportunity districts because their
incarcerated members are counted elsewhere, the apparent minority voting-age
population (MVAP) of those districts is reduced. In borderline districts---
those close to the 50\% minority threshold that defines a majority-minority
district under the \textit{Gingles} framework---this reduction can affect
whether the district is legally required to be drawn as a majority-minority
district at all.
